\section{Global form of the Standard Model gauge group}
\label{sec:global-form}

The local Standard Model gauge algebra does not by itself fix the global gauge group.  The global form may be written as
\begin{equation}
  \GSM=\frac{SU(3)_C\times SU(2)_W\times U(1)_Y}{\Gamma},
  \qquad
  \Gamma\subset \mathbb Z_6.
  \label{eq:gsm-global}
\end{equation}
The possibilities \(\Gamma=1,\mathbb Z_2,\mathbb Z_3,\mathbb Z_6\) define different spectra of line operators and different topological sectors.  Local correlation functions in flat spacetime do not fix this choice, while line operators and global probes can distinguish it \cite{TongLineOperators2017,TongGauge}.

This global-form issue is relevant to any topological or compactification completion of the DeVries construction.  A model that derives the electroweak secular equation should also specify which line operators exist and which quotient \(\Gamma\) is realized.  Recent work proposes topological transport coefficients that can distinguish the global form in analogy with fractional quantum Hall response \cite{HsinGomis2024}.  Such tests are not inputs to the mass ratio, but they constrain the class of completions.

\paragraph*{Section status.}  The global-form appendix is source-supported and functions as a completion constraint.
